\section{Model and Estimation}
\label{sec:method}

\subsection{Base-model selection as a discrete choice}
Following \citet{overgoor2019choosing}, each first adoption $k$ is a choice of one parent $j_k$ from the market
$\M_k$. Developer $k$ values candidate $j$ at $U_{kj} = \theta^\top x_{kj} + \varepsilon_{kj}$ with i.i.d.\ Gumbel
errors, which gives the conditional logit \citep{mcfadden1974conditional}
\begin{equation}
  P(j \mid \M_k) = \frac{\exp(\theta^\top x_{kj})}{\sum_{l \in \M_k} \exp(\theta^\top x_{kl})}.
  \label{eq:logit}
\end{equation}
The features $x_{kj}$ are the four groups of \cref{sec:data}; they vary over candidates and over time and, for
demand match and inertia, over developers. With $\theta = 0$ the model is \emph{uniform attachment}, every
alive model equally likely; preferential attachment is the special case in which only $\log d$ and $\1[d>0]$
enter. We interpret $\theta$ descriptively, as the weights that best summarize which signals choices co-move with,
not as causal effects (\cref{sec:discussion}).

\subsection{Sampling the competitors that matter}
\label{sec:sampling}
Evaluating \eqref{eq:logit} over a median of \marketMedian{} alternatives for \nFirst{} choices is infeasible, so we
sample a candidate set $\D_k \subset \M_k$. Each alternative enters independently (Poisson sampling) with
\begin{equation}
  \pi_{kl} = \min\{1,\; c_k (d_{kl} + 1)\}, \qquad \textstyle\sum_{l \in \M_k} \pi_{kl} = S,
  \label{eq:pps}
\end{equation}
where $d_{kl}$ is the degree of $l$ just before choice $k$ and $c_k$ is solved in closed form from the degree
histogram of the market; the chosen parent is always included. We use $S=300$; on average 8.6 candidates per choice
have $\pi=1$, so the incumbents a developer is most likely to weigh are present in every set, and the rest of the
set is a probability-proportional-to-size sample of the long tail.

\begin{proposition}[Sampling correction]
\label{prop:pps}
Under \eqref{eq:logit} and the sampling scheme \eqref{eq:pps},
\begin{equation}
  P(j \mid \D_k) = \frac{\exp(\theta^\top x_{kj} - \log \pi_{kj})}{\sum_{l \in \D_k} \exp(\theta^\top x_{kl} - \log
  \pi_{kl})},
  \label{eq:corrected}
\end{equation}
and maximizing $\sum_k \log P(j_k \mid \D_k)$ gives a consistent estimator of $\theta$.
\end{proposition}
The proof (\cref{app:sampling}) applies McFadden's positive-conditioning argument \citep{mcfadden1978modelling}:
given that $j$ was chosen, the probability of drawing $\D$ is $\prod_{l \in \D \setminus j} \pi_l
\prod_{l \notin \D}(1-\pi_l)$, so the ratio for two members of $\D$ is $\pi_l / \pi_j$. The offset $-\log \pi$
makes uniform attachment, $\theta = 0$, the natural baseline: it assigns every sampled candidate the weight
$1/\pi$, i.e.\ every model in the full market the same probability.

Why sample proportionally to degree? Consistency holds for any $\pi > 0$, but efficiency does not. A uniform sample
of 300 from a market of $3 \times 10^5$ almost never contains the dozen models that receive most adoptions, so each
choice compares the chosen parent mostly with never-used models, and the likelihood learns little about how
developers trade off one popular model against another. \Cref{sec:robustq} shows that this is not only a matter of
precision: with uniform sets the estimated weights change.

\subsection{Estimation and inference}
The training data have 35.2 million (choice, candidate) rows. We group choices into buckets of similar set size,
keep the 26 features on the GPU in single precision and accumulate the log-likelihood, gradient and Hessian in double
precision; L-BFGS with an analytic Hessian fits the full model in under a minute on one A100. Standard errors are
clustered by author \citep{cameron2015practitioner}: \nFirstAuthors{} authors make the \nFirst{} first adoptions, and
the choices of one author are strongly dependent (\cref{app:validation}). For every derived quantity, including the
Shapley shares below, we use a multiplier (Bayesian) bootstrap over authors \citep{rubin1981bayesian}: each
replicate draws a weight $w_a \sim \mathrm{Exp}(1)$ per author, reweights the author's choices and refits from the
point estimate.

\subsection{Decomposing the likelihood gain}
\label{sec:shapley}
The four groups are correlated (popular models also tend to be well documented and to match common tasks), so the
change in fit from adding a group depends on which groups are already in the model. We therefore attribute the
log-likelihood gain over uniform attachment, $G(\mathcal{S}) = \ell(\hat\theta_{\mathcal{S}}) - \ell(0)$ for the
model with the groups in $\mathcal{S}$, by the Shapley value \citep{shapley1953value}
\begin{equation}
  \phi_g = \sum_{\mathcal{S} \subseteq \mathcal{G} \setminus \{g\}}
  \frac{|\mathcal{S}|!\,(|\mathcal{G}|-|\mathcal{S}|-1)!}{|\mathcal{G}|!}
  \big[G(\mathcal{S} \cup \{g\}) - G(\mathcal{S})\big],
  \label{eq:shapley}
\end{equation}
and report shares $\phi_g / G(\mathcal{G})$. With four groups this needs $2^4 = 16$ fits, each warm-started from
a subset one group smaller. We compute the shares both in-sample and on the holdout (mean per-choice
log-likelihood), where a group can only gain by predicting future choices.

\subsection{Bounding general quality}
\label{sec:bounds}
Benchmark scores cover few models, so a small quality share could reflect missing measurement rather than
indifference to quality. We bound the quality share from above by adding an intercept $\phi_j$ for each of the
500 most adopted parents (assigned ex ante, \cref{sec:data}). The fixed effects absorb every time-invariant
attribute of a hub: unmeasured quality, but also brand, tutorials and default status. Entering the fixed effects as
a fifth Shapley player and counting their entire share as quality gives the largest share that any fixed hub
attribute could contribute. To see how much of this bound is quality, we regress $\hat\phi_j$ on external quality
measures (Open LLM Leaderboard v2 average, arena ratings; controls for parameter count and re-distributor accounts)
and correct the $R^2$ for estimation noise using the bootstrap variance of $\hat\phi_j$; the noise-corrected share
of $\mathrm{Var}(\phi)$ explained by quality splits the fixed-effect share into a quality and a popularity part.

\subsection{Does the decomposition recover known mechanisms?}
\label{sec:validation}
Two checks support the procedure (details in \cref{app:validation}). First, on synthetic preferential-attachment
networks the corrected likelihood recovers the attachment exponent under all three sampling designs. Second, we
replace every observed choice with a draw from a known utility on the same candidate sets and re-run the full
decomposition. When choices are drawn from the popularity terms alone, the decomposition attributes
\recSfPop\% of the gain to popularity; when drawn from quality and demand match alone, it attributes \recQfFit\% to
fit; when drawn from the full estimated model, it returns \recFullPop\%, \recFullMatch\% and \recFullInertia\% for
popularity, demand match and inertia, the shares estimated on the real choices. The decomposition can tell the
mechanisms apart, and it does not manufacture a popularity share where none was generated.
